% GENERATED FILE. Edit canonical lessons, then run npm run generate:books.
% canonical-source-hash: fnv1a64:67099dfc56761630
% canonical-lessons: BN-C345-songsod, BN-C345-songbidhan, BN-C345-prostab, BN-C345-onumodon-kora, BN-C345-batil-kora, BN-R345-review-338-345

\chapter{Making Laws}
\label{ch:bn-making-laws}
\hlchaptermodality{345}

\begin{chapteropening}
\textbf{By the end of this chapter:} \emph{I can say how a law is proposed, passed or thrown out.}

Complete the last lesson of chapter 345: I can say how a law is proposed, passed or thrown out.
\end{chapteropening}

\section[\texorpdfstring{\bn{সংসদ} (sôngsôd)}{সংসদ (sôngsôd)}]{\bn{সংসদ} (sôngsôd) — a parliament}

\label{lesson:BN-C345-songsod}



\begin{quote}
Before the new one: say the Bengali for power, then the Bengali for the people.
\end{quote}

\subsection*{You'll want to know: \bn{সংসদ}}
\textbf{\bn{সংসদ}} — \emph{sôngsôd} — \textquotedblleft{}a parliament\textquotedblright{}.

\begin{quote}
\bn{সংসদে} \bn{আজ} \bn{নতুন} \bn{আইন} \bn{নিয়ে} \bn{আলোচনা} \bn{হল।}
\end{quote}

\begin{quote}
\emph{sôngsôde āj nôtun āin niye ālochônā hôlô.}
\end{quote}

\begin{quote}
Parliament discussed the new law today.
\end{quote}

\begin{cousinweb}
\textbf{\bn{সংসদ}} (\emph{sôngsôd}) is Sanskrit \emph{saṃsad}, 'sitting together', an assembly (Monier-Williams). Modern Bengali uses it for a national parliament.
\end{cousinweb}

\begin{grammarlens}[title={What you've built}]
One more word for giving reasons and opinions.
\end{grammarlens}

\subsection*{Your turn}
\begin{itemize}
\raggedright
  \item \emph{Say it:} \emph{sôngsôd}
  \item \emph{Say it:} \emph{sôngsôd}, once more
  \item \emph{Say it:} \emph{sôngsôde āj nôtun āin niye ālochônā hôlô.} — Parliament discussed the new law today.
  \item \emph{Recall:} say the Bengali for power, then the Bengali for the people, then say \emph{sôngsôd} again
\end{itemize}

\begin{tcolorbox}[breakable,colback=teal!4,colframe=teal!35!black,title={Before you move on}]
What does \bn{সংসদ} mean? (\textquotedblleft{}A parliament\textquotedblright{}.) Say it once more.
\end{tcolorbox}
\section[\texorpdfstring{\bn{সংবিধান} (sôngbidhān)}{সংবিধান (sôngbidhān)}]{\bn{সংবিধান} (sôngbidhān) — a constitution}

\label{lesson:BN-C345-songbidhan}



\begin{quote}
Before the new one: say the Bengali for the people, then the Bengali for a parliament.
\end{quote}

\subsection*{You'll want to know: \bn{সংবিধান}}
\textbf{\bn{সংবিধান}} — \emph{sôngbidhān} — \textquotedblleft{}a constitution\textquotedblright{}.

\begin{quote}
\bn{সংবিধান} \bn{সব} \bn{নাগরিককে} \bn{সমান} \bn{অধিকার} \bn{দেয়।}
\end{quote}

\begin{quote}
\emph{sôngbidhān sôb nāgôrikke sômān ôdhikār dey.}
\end{quote}

\begin{quote}
The constitution gives every citizen equal rights.
\end{quote}

\begin{cousinweb}
\textbf{\bn{সংবিধান}} (\emph{sôngbidhān}) is Sanskrit \emph{saṃvidhāna}, an arrangement (Monier-Williams). The sense 'a constitution', the arrangement a state is built on, is a modern coinage.
\end{cousinweb}

\begin{grammarlens}[title={What you've built}]
One more word for giving reasons and opinions.
\end{grammarlens}

\subsection*{Your turn}
\begin{itemize}
\raggedright
  \item \emph{Say it:} \emph{sôngbidhān}
  \item \emph{Say it:} \emph{sôngbidhān}, once more
  \item \emph{Say it:} \emph{sôngbidhān sôb nāgôrikke sômān ôdhikār dey.} — The constitution gives every citizen equal rights.
  \item \emph{Recall:} say the Bengali for the people, then the Bengali for a parliament, then say \emph{sôngbidhān} again
\end{itemize}

\begin{tcolorbox}[breakable,colback=teal!4,colframe=teal!35!black,title={Before you move on}]
What does \bn{সংবিধান} mean? (\textquotedblleft{}A constitution\textquotedblright{}.) Say it once more.
\end{tcolorbox}
\section[\texorpdfstring{\bn{প্রস্তাব} (prôstāb)}{প্রস্তাব (prôstāb)}]{\bn{প্রস্তাব} (prôstāb) — a proposal; a motion}

\label{lesson:BN-C345-prostab}



\begin{quote}
Before the new one: say the Bengali for a parliament, then the Bengali for a constitution.
\end{quote}

\subsection*{You'll want to know: \bn{প্রস্তাব}}
\textbf{\bn{প্রস্তাব}} — \emph{prôstāb} — \textquotedblleft{}a proposal; a motion\textquotedblright{}.

\begin{quote}
\bn{আমার} \bn{একটা} \bn{প্রস্তাব} \bn{আছে।}
\end{quote}

\begin{quote}
\emph{āmār ekṭā prôstāb āchhe.}
\end{quote}

\begin{quote}
I have a proposal.
\end{quote}

\begin{cousinweb}
\textbf{\bn{প্রস্তাব}} (\emph{prôstāb}) is Sanskrit \emph{prastāva}, an introduction or prelude (Monier-Williams). Bengali uses it for what is put forward first in a meeting, a proposal or a motion; that use is a later one.
\end{cousinweb}

\begin{grammarlens}[title={What you've built}]
One more word for giving reasons and opinions.
\end{grammarlens}

\subsection*{Your turn}
\begin{itemize}
\raggedright
  \item \emph{Say it:} \emph{prôstāb}
  \item \emph{Say it:} \emph{prôstāb}, once more
  \item \emph{Say it:} \emph{āmār ekṭā prôstāb āchhe.} — I have a proposal.
  \item \emph{Recall:} say the Bengali for a parliament, then the Bengali for a constitution, then say \emph{prôstāb} again
\end{itemize}

\begin{tcolorbox}[breakable,colback=teal!4,colframe=teal!35!black,title={Before you move on}]
What does \bn{প্রস্তাব} mean? (\textquotedblleft{}A proposal; a motion\textquotedblright{}.) Say it once more.
\end{tcolorbox}
\section[\texorpdfstring{\bn{অনুমোদন} \bn{করা} (ônumodôn kôrā)}{অনুমোদন করা (ônumodôn kôrā)}]{\bn{অনুমোদন} \bn{করা} (ônumodôn kôrā) — to approve; to pass (a bill)}

\label{lesson:BN-C345-onumodon-kora}



\begin{quote}
Before the new one: say the Bengali for a constitution, then the Bengali for a proposal.
\end{quote}

\subsection*{You'll want to know: \bn{অনুমোদন} \bn{করা}}
\textbf{\bn{অনুমোদন} \bn{করা}} — \emph{ônumodôn kôrā} — \textquotedblleft{}to approve; to pass (a bill)\textquotedblright{}.

\begin{quote}
\bn{সংসদ} \bn{প্রস্তাবটা} \bn{অনুমোদন} \bn{করল।}
\end{quote}

\begin{quote}
\emph{sôngsôd prôstābṭā ônumodôn kôrlô.}
\end{quote}

\begin{quote}
Parliament approved the proposal.
\end{quote}

Past (he/she …): \bn{তিনি} \bn{পরিকল্পনাটা} \bn{অনুমোদন} \bn{করলেন}, \emph{tini pôrikôlpônāṭā ônumodôn kôrlen}.

\begin{cousinweb}
\textbf{\bn{অনুমোদন} \bn{করা}} (\emph{ônumodôn kôrā}) is Sanskrit \emph{anumodana}, assent or applause (Monier-Williams), with \bn{করা}: to give formal approval.
\end{cousinweb}

\begin{grammarlens}[title={What you've built}]
One more word for giving reasons and opinions.
\end{grammarlens}

\subsection*{Your turn}
\begin{itemize}
\raggedright
  \item \emph{Say it:} \emph{ônumodôn kôrā}
  \item \emph{Say it:} \emph{ônumodôn kôrā}, once more
  \item \emph{Say it:} \emph{sôngsôd prôstābṭā ônumodôn kôrlô.} — Parliament approved the proposal.
  \item \emph{Recall:} say the Bengali for a constitution, then the Bengali for a proposal, then say \emph{ônumodôn kôrā} again
\end{itemize}

\begin{tcolorbox}[breakable,colback=teal!4,colframe=teal!35!black,title={Before you move on}]
What does \bn{অনুমোদন} \bn{করা} mean? (\textquotedblleft{}To approve; to pass (a bill)\textquotedblright{}.) Say it once more.
\end{tcolorbox}
\section[\texorpdfstring{\bn{বাতিল} \bn{করা} (bātil kôrā)}{বাতিল করা (bātil kôrā)}]{\bn{বাতিল} \bn{করা} (bātil kôrā) — to cancel; to repeal}

\label{lesson:BN-C345-batil-kora}



\begin{quote}
Before the new one: say the Bengali for a proposal, then the Bengali for to approve.
\end{quote}

\subsection*{You'll want to know: \bn{বাতিল} \bn{করা}}
\textbf{\bn{বাতিল} \bn{করা}} — \emph{bātil kôrā} — \textquotedblleft{}to cancel; to repeal\textquotedblright{}.

\begin{quote}
\bn{সরকার} \bn{পুরনো} \bn{আইনটা} \bn{বাতিল} \bn{করল।}
\end{quote}

\begin{quote}
\emph{sôrkār purôno āinṭā bātil kôrlô.}
\end{quote}

\begin{quote}
The government repealed the old law.
\end{quote}

Past (he/she …): \bn{তিনি} \bn{মিটিংটা} \bn{বাতিল} \bn{করলেন}, \emph{tini miṭingṭā bātil kôrlen}.

\begin{cousinweb}
\textbf{\bn{বাতিল} \bn{করা}} (\emph{bātil kôrā}) is usually traced to Arabic \emph{bāṭil}, void, through Persian, with \bn{করা}: to make a law, a ticket or a meeting void.
\end{cousinweb}

\begin{grammarlens}[title={What you've built}]
One more word for giving reasons and opinions.
\end{grammarlens}

\subsection*{Your turn}
\begin{itemize}
\raggedright
  \item \emph{Say it:} \emph{bātil kôrā}
  \item \emph{Say it:} \emph{bātil kôrā}, once more
  \item \emph{Say it:} \emph{sôrkār purôno āinṭā bātil kôrlô.} — The government repealed the old law.
  \item \emph{Recall:} say the Bengali for a proposal, then the Bengali for to approve, then say \emph{bātil kôrā} again
\end{itemize}

\begin{tcolorbox}[breakable,colback=teal!4,colframe=teal!35!black,title={Before you move on}]
What does \bn{বাতিল} \bn{করা} mean? (\textquotedblleft{}To cancel; to repeal\textquotedblright{}.) Say it once more.
\end{tcolorbox}
\section[(dialogue)]{Looking back: from Documents and Facts to Making Laws}

\label{lesson:BN-R345-review-338-345}



\begin{quote}
Say the Bengali for a certificate and information.
\end{quote}

\subsection*{Your turn}
Say these aloud:

\begin{itemize}
\raggedright
  \item a certificate, information, testimony, to check against, to refute
  \item a right, a responsibility, self-interest, to defend, to enjoy a right
  \item lawful, banned, a clause, to carry out, to enforce
  \item to violate, to disobey, forgery, corruption, an injustice
  \item a lawyer, to try a case, a verdict, a formal complaint, to settle a dispute
  \item blame, responsible for a fault, compensation, to answer for, to let off
  \item a nation-state, to run an organisation, an administration, power, the people
  \item a parliament, a constitution, a proposal, to approve, to repeal
\end{itemize}

\begin{tcolorbox}[breakable,colback=teal!4,colframe=teal!35!black,title={Before you move on}]
A certificate? (\textbf{\bn{সনদ}}, \emph{sônôd}.) Information? (\textbf{\bn{তথ্য}}, \emph{tôthyô}.) Testimony? (\textbf{\bn{সাক্ষ্য}}, \emph{sākkhyô}.) To check against? (\textbf{\bn{মেলানো}}, \emph{melāno}.) To refute? (\textbf{\bn{খণ্ডন} \bn{করা}}, \emph{khônḍôn kôrā}.) A right? (\textbf{\bn{অধিকার}}, \emph{ôdhikār}.) A responsibility? (\textbf{\bn{দায়িত্ব}}, \emph{dāyitvô}.) Self-interest? (\textbf{\bn{স্বার্থ}}, \emph{svārthô}.) To defend? (\textbf{\bn{রক্ষা} \bn{করা}}, \emph{rôkkhā kôrā}.) To enjoy a right? (\textbf{\bn{ভোগ} \bn{করা}}, \emph{bhog kôrā}.) Lawful? (\textbf{\bn{বৈধ}}, \emph{boidhô}.) Banned? (\textbf{\bn{নিষিদ্ধ}}, \emph{nishiddhô}.) A clause? (\textbf{\bn{ধারা}}, \emph{dhārā}.) To carry out? (\textbf{\bn{পালন} \bn{করা}}, \emph{pālôn kôrā}.) To enforce? (\textbf{\bn{প্রয়োগ} \bn{করা}}, \emph{prôyog kôrā}.) To violate? (\textbf{\bn{লঙ্ঘন} \bn{করা}}, \emph{lôngghôn kôrā}.) To disobey? (\textbf{\bn{অমান্য} \bn{করা}}, \emph{ômānyô kôrā}.) Forgery? (\textbf{\bn{জালিয়াতি}}, \emph{jāliyāti}.) Corruption? (\textbf{\bn{দুর্নীতি}}, \emph{durnīti}.) An injustice? (\textbf{\bn{অন্যায়}}, \emph{ônyāy}.) A lawyer? (\textbf{\bn{উকিল}}, \emph{ukil}.) To try a case? (\textbf{\bn{বিচার} \bn{করা}}, \emph{bichār kôrā}.) A verdict? (\textbf{\bn{রায়}}, \emph{rāy}.) A formal complaint? (\textbf{\bn{নালিশ}}, \emph{nālish}.) To settle a dispute? (\textbf{\bn{মীমাংসা} \bn{করা}}, \emph{mīmāngsā kôrā}.) Blame? (\textbf{\bn{দোষ}}, \emph{dosh}.) Responsible for a fault? (\textbf{\bn{দায়ী}}, \emph{dāyī}.) Compensation? (\textbf{\bn{ক্ষতিপূরণ}}, \emph{khôtipūrôn}.) To answer for? (\textbf{\bn{জবাবদিহি} \bn{করা}}, \emph{jôbābdihi kôrā}.) To let off? (\textbf{\bn{রেহাই} \bn{দেওয়া}}, \emph{rehāi deoyā}.) A nation-state? (\textbf{\bn{রাষ্ট্র}}, \emph{rāshṭrô}.) To run an organisation? (\textbf{\bn{পরিচালনা} \bn{করা}}, \emph{pôrichālônā kôrā}.) An administration? (\textbf{\bn{প্রশাসন}}, \emph{prôshāsôn}.) Power? (\textbf{\bn{ক্ষমতা}}, \emph{khômôtā}.) The people? (\textbf{\bn{জনগণ}}, \emph{jônôgôn}.) A parliament? (\textbf{\bn{সংসদ}}, \emph{sôngsôd}.) A constitution? (\textbf{\bn{সংবিধান}}, \emph{sôngbidhān}.) A proposal? (\textbf{\bn{প্রস্তাব}}, \emph{prôstāb}.) To approve? (\textbf{\bn{অনুমোদন} \bn{করা}}, \emph{ônumodôn kôrā}.) To repeal? (\textbf{\bn{বাতিল} \bn{করা}}, \emph{bātil kôrā}.)
\end{tcolorbox}
